\documentclass[12pt,a4paper]{article}
\usepackage[utf8]{inputenc}
\usepackage[vietnamese]{babel}
\usepackage[margin=2cm]{geometry}
\usepackage{tcolorbox}
\usepackage{enumitem}

\pagestyle{plain}

\begin{document}

% --- HEADER 2 CỘT NGHỊ ĐỊNH 30/2020/NĐ-CP ---
\noindent
\begin{minipage}[t]{0.45\textwidth}
    \begin{center}
        \textbf{CÔNG AN THÀNH PHỐ}\\
        \textbf{\underline{PHÒNG CẢNH SÁT HÌNH SỰ}}\\[0.2cm]
        \small{Số: 02/BB-KNHT}\\
        \textit{V/v khám nghiệm hiện trường vụ án}
    \end{center}
\end{minipage}
\hfill
\begin{minipage}[t]{0.52\textwidth}
    \begin{center}
        \textbf{CỘNG HÒA XÃ HỘI CHỦ NGHĨA VIỆT NAM}\\
        \textbf{\underline{Độc lập - Tự do - Hạnh phúc}}\\[0.2cm]
        \textit{Hà Nội, ngày 25 tháng 07 năm 2026}
    \end{center}
\end{minipage}

\vspace{0.8cm}

% --- TIÊU ĐỀ VĂN BẢN ---
\begin{center}
    \textbf{\LARGE BIÊN BẢN KHÁM NGHIỆM HIỆN TRƯỜNG}\\
    \textbf{\large (CHUYÊN ÁN TEST-99)}
\end{center}

\vspace{0.4cm}

\begin{tcolorbox}[colback=gray!8!white,colframe=black!75!black,title=\textbf{THÔNG TIN THỜI GIAN \& ĐỊA ĐIỂM}]
\begin{itemize}[leftmargin=*]
    \item \textbf{Địa điểm:} Căn nhà số 14, Đường Bờ Sông, Khu giải tỏa Phân khu Cảng phía Bắc.
    \item \textbf{Thời gian tiến hành:} 08:30 đến 11:30 ngày 25/07/2026.
    \item \textbf{Chủ trì khám nghiệm:} Đại úy Lê Minh - Trưởng ban chuyên án.
\end{itemize}
\end{tcolorbox}

\section*{I. MÔ TẢ TỔNG QUAN HIỆN TRƯỜNG}
\begin{itemize}
    \item Căn nhà 1 tầng cũ diện tích khoảng 75m2, nằm trong khu vực giải tỏa đền bù hẻo lánh. Căn nhà không trang bị hệ thống camera an ninh.
    \item Cửa chính khóa xích ngoắc lỏng; cửa sau (đi ra hẻm nhỏ hướng ra sông) bị bật chốt trong, có dấu vết bị đẩy mở từ bên ngoài.
\end{itemize}

\section*{II. CHI TIẾT TẠI PHÒNG KHÁCH (NƠI PHÁT HIỆN THI THỂ)}

\subsection*{1. Vị trí thi thể \& Vũng máu}
\begin{itemize}
    \item Nạn nhân nằm ngửa trên sàn gạch cũ cạnh bàn trà gỗ.
    \item Vũng máu đầm đìa lan từ vùng cổ nạn nhân ra xung quanh.
\end{itemize}

\subsection*{2. Dấu vết xáo trộn \& Vật chứng rơi vỡ}
\begin{itemize}
    \item Bộ bình cốc thủy tinh uống trà bị đập vỡ vụn dưới sàn nhà gần mép bàn.
    \item Tìm thấy 01 mảnh thủy tinh vỡ sắc nhọn dài khoảng 8cm dính máu khô và có dấu ấn vân tay dạng miết sát cạnh thi thể.
    \item Tủ âm tường gỗ ở góc phòng bị xê dịch. Khóa khoang tủ sắt bên trong bị đập bung.
    \item Nhiều tờ bản vẽ địa chính, đơn từ đền bù đất đai và một mảnh giấy di chúc viết tay bị văng ra sàn nhà gần khoang tủ âm tường.
\end{itemize}

\section*{III. DẤU VẾT THU THỨC ĐƯỢC}
\begin{enumerate}
    \item \textbf{Vật chứng 01:} 01 Mảnh thủy tinh sắc nhọn dính máu khô (\texttt{p3}).
    \item \textbf{Vật chứng 02:} Mẫu máu và dấu giày nam size 41 gần khu vực bàn trà.
    \item \textbf{Vật chứng 03:} Các tập hồ sơ giấy tờ vương vãi bao gồm đơn đăng ký đền bù và bản sao di chúc.
\end{enumerate}

\vspace{1cm}

% --- CHỮ KÝ VÀ NƠI NHẬN CHUẨN NGHỊ ĐỊNH 30 ---
\noindent
\begin{minipage}[t]{0.45\textwidth}
    \textbf{\textit{Nơi nhận:}}\\
    \small{- Thủ trưởng Cơ quan CSĐT (b/c);}\\
    \small{- Viện KSMND Thành phố (để k/s);}\\
    \small{- Đội Cảnh sát Điều tra án mạng;}\\
    \small{- Lưu: VT, CSHS.}
\end{minipage}
\hfill
\begin{minipage}[t]{0.5\textwidth}
    \begin{center}
        \textbf{TRƯỞNG BAN CHUYÊN ÁN}\\
        \textit{(Ký, ghi rõ họ tên và đóng dấu)}\\[1.5cm]
        \textbf{Đại úy Lê Minh}
    \end{center}
\end{minipage}

\end{document}
